\documentclass[12pt]{article}
\usepackage[utf8]{inputenc}
\usepackage[T1]{fontenc}
\usepackage[margin=1in]{geometry}
\usepackage{setspace}
\usepackage{xcolor}
\usepackage[colorlinks=true, urlcolor=blue!60!black, linkcolor=blue!60!black]{hyperref}
\setlength{\parindent}{2em}
\setlength{\parskip}{1.25em}
\renewcommand{\baselinestretch}{1.0}
\begin{document}
\begin{center}
{\LARGE \textit{Planned Parenthood of Southeastern Pennsylvania v.\ Casey}, 505 U.S.\ 833 (1992)\par}
\vspace{0.8em}
{\large Case Brief --- Four-Page Edition\par}
\end{center}
\vspace{1em}

\section{Memory Jogger}

In a joint opinion by Justices O'Connor, Kennedy, and Souter, the Court reaffirmed the ``essential holding'' of \textit{Roe v.\ Wade} on stare decisis and liberty grounds but replaced the trimester framework with the ``undue burden'' standard, upholding four Pennsylvania abortion regulations while striking spousal notification; \textit{Dobbs} (2022) later overruled it (\href{https://supreme.justia.com/cases/federal/us/505/833/}{505 U.S.\ 833}).

\section{Detailed Case Facts}

Amendments to Pennsylvania's Abortion Control Act required (1) informed consent with a 24-hour waiting period, including information about the ``probable gestational age of the unborn child'' and state materials on fetal development and adoption; (2) one parent's consent for minors, with a judicial bypass; (3) a married woman's signed statement that she had notified her husband, with narrow exceptions; (4) a medical-emergency exception; and (5) recordkeeping and reporting, including a married woman's reason for not notifying her husband. Five clinics led by Planned Parenthood and a physician class sued Governor Robert Casey before the provisions took effect, alleging all five were facially unconstitutional.

\section{Procedural History}

After a bench trial, the district court struck all five provisions under \textit{Roe}. The Third Circuit, reading \textit{Webster} and \textit{Hodgson} to adopt an undue-burden test, upheld all but spousal notification, over then-Judge Alito's partial dissent. The United States, as \textit{amicus}, urged the Court to overrule \textit{Roe}. The Court decided the case on June 29, 1992.

\section{Judicial Votes}

The Court fractured across five opinions.
\begin{itemize}
\item \textbf{Joint opinion} (O'Connor, Kennedy, Souter): Opinion of the Court in Parts I, II, III, V-A, V-C, and VI (with Blackmun and Stevens); plurality in Parts IV, V-B, and V-D.
\item \textbf{Stevens} and \textbf{Blackmun:} each concurred in part and dissented in part.
\item \textbf{Rehnquist} and \textbf{Scalia}, each joined by White, Thomas, and the other: concurred in the judgment in part and dissented in part.
\item \textbf{Tallies:} reaffirming \textit{Roe}, 5--4; informed consent and waiting period, 7--2 upheld; parental consent, 7--2 upheld; spousal notification, 5--4 struck; general reporting, 8--1 upheld.
\end{itemize}

\section{Holding}

A five-Justice majority reaffirmed \textit{Roe}'s essential holding: a right to choose abortion before viability without undue interference; state power to restrict post-viability abortion with life and health exceptions; and legitimate state interests from the outset in maternal health and potential life. The controlling plurality replaced the trimester framework and strict scrutiny with the undue burden standard, judging facial challenges by the ``large fraction'' of women for whom a law is an actual restriction. Spousal notification and its reporting requirement were invalid; the other provisions were upheld, and \textit{Akron I} and \textit{Thornburgh} were overruled in part.

\section{Analysis of Opinions}

\textbf{Joint opinion.} Liberty under the Due Process Clause is a ``rational continuum'' (Harlan in \textit{Poe v.\ Ullman}) linked to \textit{Griswold}, \textit{Eisenstadt}, and \textit{Loving}. Under stare decisis, \textit{Roe} was workable, women had relied on it, no doctrine had eroded it, and earlier viability did not undermine the line. Unlike \textit{West Coast Hotel} and \textit{Brown}, overruling ``under fire'' would cause ``profound and unnecessary damage to the Court's legitimacy.'' The plurality found the trimester framework undervalued the interest in potential life. Spousal notification, in light of domestic-violence findings, was a substantial obstacle in a large fraction of relevant cases: ``[w]omen do not lose their constitutionally protected liberty when they marry.''

\textbf{Stevens.} He would bar persuasive state messaging and the waiting period, which rested on ``outmoded and unacceptable assumptions about the decisionmaking capacity of women.''

\textbf{Blackmun.} He hailed ``an act of personal courage and constitutional principle'' but would retain strict scrutiny, grounded also in gender equality, and warned that the right survived by one vote.

\textbf{Rehnquist.} \textit{Roe} was wrong; abortion is \textit{sui generis}; the joint opinion offered a ``newly minted variation on \textit{stare decisis}''; rational basis would sustain every provision.

\textbf{Scalia.} The Constitution says nothing about abortion and tradition permitted its prohibition. He mocked ``reasoned judgment'' as a ``collection of adjectives,'' the test as a ``verbal shell game,'' and the legitimacy rationale as a ``Nietzschean vision'' of an ``Imperial Judiciary.''

\section{Examples: Future Applications}

\begin{enumerate}
\item \textbf{Same-Side: Admitting Privileges.} A law closing half the State's clinics without health benefits is a substantial obstacle (cf.\ \textit{Whole Woman's Health v.\ Hellerstedt}).
\item \textbf{Same-Side: Husband's Acknowledgment.} Requiring a husband's written acknowledgment absent a police report gives an abusive spouse a veto, like \S\ 3209.
\item \textbf{Opposite-Side: Telehealth Informed Consent.} An accurate state brochure 24 hours before an abortion, with telehealth consultation, is less burdensome than what \textit{Casey} upheld.
\item \textbf{Opposite-Side: Post-Viability Ban.} A ban after 26 weeks with life and health exceptions fits \textit{Casey}'s second core holding.
\item \textbf{Fence-Sitter: 72-Hour, Two-Trip Rule.} In a rural State with two clinics, delays average two weeks. \textit{Casey} said a ``particular burden is not of necessity a substantial obstacle'' but left open a stronger record.
\end{enumerate}

\section{Critique}

\textbf{Originalist critiques.} Following Scalia, critics find no textual or historical warrant for the right given nineteenth-century prohibitions; Michael Stokes Paulsen condemns ``super-\textit{stare decisis}'' as entrenching error; and conservatives say the undue burden test lacked objective standards.

\textbf{Progressive and feminist critiques.} Then-Judge Ruth Bader Ginsburg and Reva Siegel argued that resting the right on due process rather than equal protection left it vulnerable, though they welcomed \textit{Casey}'s anti-coverture principles. Reproductive-rights scholars warned that the malleable standard invited restrictions harming poor, rural, and minor women.

\section{Key Quotations}

\begin{quote}
``At the heart of liberty is the right to define one's own concept of existence, of meaning, of the universe, and of the mystery of human life. Beliefs about these matters could not define the attributes of personhood were they formed under compulsion of the State.''
\end{quote}
\noindent --- Joint Opinion (Opinion of the Court), 505 U.S.\ at 851.

\begin{quote}
``An undue burden exists, and therefore a provision of law is invalid, if its purpose or effect is to place a substantial obstacle in the path of a woman seeking an abortion before the fetus attains viability.''
\end{quote}
\noindent --- Joint Opinion (plurality), 505 U.S.\ at 878.

\begin{quote}
``Liberty finds no refuge in a jurisprudence of doubt.''
\end{quote}
\noindent --- Joint Opinion (Opinion of the Court), 505 U.S.\ at 844.

\end{document}
